\section{LATEX}

\LaTeX ist ein Softwarepaket, das die Benutzung des Textsatzsystems TeX mit Hilfe von Makros vereinfacht. LaTeX liegt derzeit in der Version $2\epsilon$ vor. LaTeX wurde Anfang der 1980er Jahre von Leslie Lamport entwickelt.Der Name bedeutet so viel wie Lamport TeX. Die Entwicklung wurde seit den 1990er Jahren von einer Anzahl Entwicklern weitergeführt. Heute ist LaTeX die populärste Methode, TeX zu verwenden. 
\cite{wikilatex}
%<https://de.wikipedia.org/wiki/LaTeX> 

Im Gegensatz zu anderen Textverarbeitungsprogrammen, die nach dem What-you-see-is-what-you-get-Prinzip funktionieren, arbeitet der Autor mit Textdateien, in denen er innerhalb eines Textes anders zu formatierende Passagen oder Überschriften mit Befehlen textuell auszeichnet. 
\cite{wikilatex}
%Aus <https://de.wikipedia.org/wiki/LaTeX#Grundprinzip> 

Ich benutze Latex zum Erstellen wissenschaftlicher und didaktischer Dokumente (Berichte, Paper, Doktorarbeit, Vorlesungsunterlagen) aber auch zum Erstellen meines Lebenslaufes, Bewerbungsschreiben, formale Briefe im Allgemeinen. Einmal installiert und zum Laufen gebracht, erspart es wirklich sehr viel Formatierungsärger und man kann sich beim Schreiben auf den Inhalt konzentrieren. Die Form passt automatisch.

Um mit LaTeX zu arbeiten, werden zwei Dinge benötigt: Einerseits die LaTeX-Software und andererseits eine Entwicklungsumgebung, mit der der LaTeX-Code eingegeben und die Umsetzung in ein fertig gesetztes Dokument angestoßen wird. Hier mein Paket (alles kostenfrei). Bitte in genau dieser Reihenfolge installieren.

\begin{itemize}
    \item • MiKTeX (TEX- Distribution)
\item Ghostscript (zur PS- Erstellung)
\item GSview (PS- Viewer. Ghostview eignet sich auch zur Betrachtung von pdf Dateien. Es ist in der Lage pdf Dateien zu „reloaden“ wenn sie durch einen LateX Kompiliervorgang erneuert wurden. Acrobat- Reader geht natürlich auch.)
\item JabRef, (Javaprogramm zur Verwaltung von Bibliografie-Dateien (.bib))
\item Inkscape (zum Erstellen von EPS- Grafiken)
\item TeXnicCenter (Editor)
\end{itemize}

Der Wizard legt ein paar Ausgabekonfigurationen an. Er fragt dabei meist auch nach Pfaden zu verschiedenen exe- Dateien. Sie finden sich im MikTeX- Ordner 
%(etwa da: \texttt{:\Program Files\MiKTeX 2.9\miktex\bin\x64 })
Man muss die Ausgabekonfiguration "LaTeX->PS->PDF" auswählen, damit mein Beispielprojekt funktioniert. 

HILFE ONLINE: \url{https://latex.tugraz.at/}
Dort wird das Zitieren auch sehr ausführlich erklärt und man findet Links zu den Downloads der bei der Installation benötigten Software.

OVERFEAF!!!
HIER VORLAGE!!!


